\documentclass{article}
% Source: Topic_07_Teacher_Edition.pdf, PDF pages 54-65 (printed pages 373A-380B)
\usepackage{amsmath}
\usepackage{amssymb}
\usepackage{graphicx}
\usepackage{tikz}
\usepackage{pgfplots}
\pgfplotsset{compat=1.18}
\usepackage{booktabs}
\usepackage{xcolor}

\newenvironment{lesson-meta}{}{}        % lesson code, title, objective, EQ, vocab
\newenvironment{item-analysis}{}{}      % Savvas's Example × DOK table
\newenvironment{model-discuss}{}{}      % Launch opener
\newenvironment{concept-box}{}{}        % in-lesson concept reference
\newenvironment{concept-summary}{}{}    % end-of-lesson summary
\newenvironment{example}[1]{}{}         % #1 = example number
\newenvironment{tryit}[1]{}{}           % #1 = tryit number
\newenvironment{practice}[2]{}{}        % #1 = practice number, #2 = DOK
\newenvironment{te-addendum}[2]{}{}     % #1 = anchor example, #2 = type
\newcommand{\answer}[1]{}               % answer key, inside any env
\newcommand{\placeholder}[2]{}          % #1 = type, #2 = description
\newcommand{\uncertain}[2]{#1}          % #1 = best guess, #2 = alternatives
\newcommand{\te}[1]{}                   % teacher-only inline annotation

\begin{document}
% Source page 54: 373A Lesson Overview
\begin{lesson-meta}
lesson: 7-4
title: Translating Trigonometric Functions
standards: none printed
practices: none printed
objective: I can use transformations to analyze and sketch graphs of trigonometric functions.

essential question: How can you find and use translations of graphs of trigonometric functions?

\textbf{VOCABULARY}
\item phase shift
\textbf{TOPIC VOCABULARY}
\te{No topic vocabulary list is printed on these pages.}
\begin{tabular}{ll}
Lesson Code & 7-4 \\
Title & Translating Trigonometric Functions \\
Standards & none printed \\
I CAN Objective & I can use transformations to analyze and sketch graphs of trigonometric functions. \\
Essential Question & How can you find and use translations of graphs of trigonometric functions? \\
Lesson Vocabulary & phase shift \\
Topic Vocabulary & none printed \\
\end{tabular}
\end{lesson-meta}
\begin{te-addendum}{0}{GeneralizeETP}
\textbf{Lesson Overview: FOCUS}
Objective. Students will be able to:
\item Identify how changing the parameters of the sine or cosine function affects the graph of the function.
\item Use trigonometric functions to model situations with specified amplitude, frequency, and midline.
Essential Understanding. The parameters (a), (b), (c), and (d) represent values that transform the functions (y=a\sin b(x-c)+d) and (y=a\cos b(x-c)+d). (|a|) represents the amplitude of the function. The period of the function is (\frac{2\pi}b). (c) represents a horizontal shift, called a phase shift; and (d) represents a vertical shift
\textbf{COHERENCE}
Previously in this topic students:
\item Identified the effects of vertical and horizontal stretches and compressions on the graphs of cosine and sine functions.
\item Graphed functions of the form (y=a\sin bx) and (y=a\cos bx).
In this lesson, students:
\item Graph phase and vertical shifts of functions of the form (y=a\sin b(x-c)+d) and (y=a\cos b(x-c)+d).
Later in this topic students will:
\item Graph the tangent and cotangent functions, as well as the secant and cosecant functions.
\textbf{RIGOR}
This lesson emphasizes a blend of conceptual understanding and application.
\item Students understand how the values of (a), (b), (c), and (d) affect each feature of the function (y=a\sin b(x-c)+d).
\item Students use the function in the form (y=a\sin b(x-c)+d) to model real-world situations, such as the temperature fluctuations in a city.
SavvasRealize.com. Program Resources.
\end{te-addendum}
\begin{te-addendum}{0}{ELLAddendum}
Glossary is available at SavvasRealize.com.
\textbf{Vocabulary Builder}
\begin{tabular}{ll}
English & Spanish \\
REVIEW VOCABULARY & \\
transformation & transformación \\
translation & traslación \\
NEW VOCABULARY & \\
phase shift & cambio de fase \\
\end{tabular}
\textbf{VOCABULARY ACTIVITY}
Sketch the graph of the quadratic function (f(x)=x^2) and the graphs of both a horizontal and a vertical translation of (f). Write the equations of each translated function and discuss the relationship between the equations and the translations. Review the other types of transformations. Then introduce the new term phase shift. Have students complete the following sentences.
1. A change to the equation of a function resulting in dilation, reflection, rotation, or translation is a \underline{\hspace{1cm}}.
\answer{transformation}
2. The shift of a function vertically, horizontally, or both without changing its shape or orientation is a \underline{\hspace{1cm}}.
\answer{translation}
3. A horizontal translation of a periodic function is a \underline{\hspace{1cm}}.
\answer{phase shift}
\end{te-addendum}
\begin{te-addendum}{0}{LearnTogether}
\textbf{Student Companion}
Students can do their work for the lesson in their Student Companion or on Savvas Realize.
% FIG 54 962 779 1115 958
\placeholder{illustration}{Black Student Companion book cover, purple and blue circular abstract artwork, title enVision Algebra2, SAVVAS.}
\end{te-addendum}
\begin{te-addendum}{0}{GeneralizeETP}
\textbf{Mathematics Overview}
Students graph transformation of periodic functions, such as a phase shift. They also identify the transformation when given a graph or table of data values.
\textbf{Applying Math Practices}
Look For and Make Use of Structure
Students notice that adding a constant value to (x) shifts the graph of a function horizontally and adding a constant value to (y) shifts the graph of a function vertically.
Look For and Express Regularity in Repeated Reasoning
Students make generalizations about the relationship between changes to the parameters (a), (b), (c), and (d) in the function (y=a\sin b(x-c)+d) and transformations of the graph of the function.
\end{te-addendum}
% Source page 55: 373B Explore
\begin{te-addendum}{0}{LearnTogether}
\textbf{CRITIQUE \& EXPLAIN}
INSTRUCTIONAL FOCUS Students use their knowledge of translations of functions to explore and analyze the relationship between the graphs of sine and cosine functions. Students recognize that each of these functions can also be written as a translation of the other. This prepares students to understand horizontal phase shifts in the lesson.
STUDENT COMPANION Students can complete the Critique \& Explain activity in their Student Companion.
Critique \& Explain is available at SavvasRealize.com. Activity.
\end{te-addendum}
\begin{te-addendum}{0}{GeneralizeETP}
\textbf{Before. WHOLE CLASS. Implement Tasks that Promote Reasoning and Problem Solving}
Q: What do you notice about the graph of (f(x)) compared to the graph of (g(x))?
\answer{f(x) has the same amplitude, period, and midline as g(x). f(x) is translated (\pi/2) units to the left of g(x).}
\textbf{During. SMALL GROUP. Support Productive Struggle in Learning Mathematics}
Q: How does adding or subtracting a positive constant value to the input of a function (f(x+k)) affect the output?
\answer{Adding a constant k translates the graph left k units, and subtracting a constant k translates the graph right k units.}
\textbf{For Early Finishers}
Q: How would you write an equation for a sine or cosine function that would be translated vertically? Give two examples, one function translated up and one translated down.
\answer{A function is translated vertically when a constant is added to the output of the function. For example, the graph of (h(x)=\sin(x)+3) would be translated up3 units, and the graph of (m(x)=\sin(x)-2) would be translated down2 units.}
\textbf{After. WHOLE CLASS. Facilitate Meaningful Mathematical Discourse}
Q: How would you describe the relationship between the graphs of the sine and cosine functions?
\answer{The graph of (y=\sin x) is a horizontal translation of the graph of (y=\cos x).}
\end{te-addendum}
\begin{te-addendum}{0}{HabitsOfMind}
\textbf{HABITS OF MIND. Use with CRITIQUE \& EXPLAIN}
Make Sense and Persevere Do you think there might be a similar relationship between the graphs of tangent and cotangent? Explain.
\answer{No, the graph of the tangent function is increasing everywhere, but the graph of the cotangent function is decreasing everywhere. You could not translate the graph of the tangent function and make it match the graph of the cotangent function.}
\end{te-addendum}
\begin{model-discuss}
\textbf{CRITIQUE \& EXPLAIN}
Sadie and Zhang use translations to relate the graphs of sine and cosine. Sadie claims that (\sin x=\cos(x-\frac\pi2)). Zhang insists that (\cos x=\sin(x+\frac\pi2)).
% FIG 55 1005 647 1120 712
\placeholder{graph}{Student launch: red g(x)=sin x and blue f(x)=cos x, amplitude1 period2pi, x,y arrowed axes origin O; x labels-2pi,2pi with ticks pi/2, y labels1,-1 and ticks1/2. Window roughly -5pi/2 to5pi/2,-3/2 to3/2; red upward crossing origin, blue maximum at0; curves have arrows at ends.}
A. Who is correct? How do you know?
\answer{A. Both are correct. Shifting the graph of (y=\cos x) to the right (\pi/2) units results in the graph of (y=\sin x). So, (\sin x=\cos(x-\frac\pi2)). And shifting the graph of (y=\sin x) to the left (\pi/2) units results in the graph of (y=\cos x). So, (\cos x=\sin(x+\frac\pi2)).}
B. Communicate Precisely Victor claims that knowing (\sin0=\cos\frac\pi2) can be used to determine who is correct. Is Victor's suggestion helpful? If so, explain why. If not, explain why not.
\answer{B. No; you can use this identity to show that Sadie's claim is correct, but it does not help you determine whether Zhang's claim is correct.}
\te{Printed answer treats one equality as evidence of an identity; one input does not establish either identity.}
\end{model-discuss}
\begin{te-addendum}{0}{GeneralizeETP}
\textbf{Digital CRITIQUE \& EXPLAIN. ESSENTIAL QUESTION. Go back. SOLUTION.}
Sadie and Zhang use translations to relate the graphs of sine and cosine. Sadie claims that (\sin x=\cos(x-\frac\pi2)). Zhang insists that (\cos x=\sin(x+\frac\pi2)).
A. Who is correct? How do you know? Enter your answer.
B. Communicate Precisely Victor claims that knowing (\sin\pi=\cos\frac\pi2) can be used to determine who is correct. Is Victor's suggestion helpful? If so, explain why. If not, explain why not. Enter your answer.
\te{Digital version prints sin pi; student version prints sin0.}
% FIG 55 962 246 1110 338
\placeholder{graph}{Digital sine/cosine launch graph, orange g(x)=sin x and blue f(x)=cos x. x,y axes origin O, labels x=-2pi,2pi and y=-1,1 with half-unit ticks. Both amplitude1 period2pi; blue maximum slightly right of y-axis, orange crosses upward slightly right of y-axis, although labels give unshifted functions. Curves span two periods, with blue end arrows and orange endpoints on x-axis.}
\end{te-addendum}
% Source page 56: 373 Understand & Apply
\begin{te-addendum}{0}{LearnTogether}
\textbf{INTRODUCE THE ESSENTIAL QUESTION. Establish Mathematics Goals to Focus Learning}
Students learn the general form of the sine function, (y=a\sin b(x-c)+d). Students recognize that (a) correlates with the amplitude of the function, (b) correlates with the period, and (c) and (d) affect the translation of the function.
\end{te-addendum}
\begin{te-addendum}{1}{GeneralizeETP}
\textbf{Understand Phase Shift as a Horizontal Translation. Build Procedural Fluency From Conceptual Understanding}
PART A
Q: How does the sign of (c) affect the direction of translation?
\answer{If c is positive, the graph shifts to the right. If c is negative, the graph shifts to the left.}
PART B
Q: Why do you need to add or subtract (c) from (x) rather than from (2x)?
\answer{When a graph is translated horizontally, you want to add or subtract the constant from only the input of the function, which is the x-value. The2 is the coefficient of x, and it affects the period of the function.}
Q: How is the period of the given function affected by (c)?
\answer{The period does not change.}
Example is Powered by Desmos. Activities are available at SavvasRealize.com.
\end{te-addendum}
\begin{te-addendum}{1}{PurposefulQuestions}
\textbf{Additional Examples. ADDITIONAL EXAMPLE 1. ESSENTIAL QUESTION. Go back. SOLUTION.}
Compare the phase shifts of the following functions:
(f(x)=\frac12\sin(x-\frac\pi4)) and (g(x)=\frac12\sin(x+\pi))?
\answer{(h=\frac\pi4) and (h=-\pi); This means that function f is shifted right (\pi/4) units and function g is shifted left (\pi) units.}
Example1 Help students determine how the value of (h) affects the phase shift of a sine function.
Q: How does the magnitude of the value for (h) affect the horizontal translation or phase shift of a sine function?
\answer{The larger the value for h, the farther the function is shifted horizontally.}
\end{te-addendum}
\begin{te-addendum}{5}{PurposefulQuestions}
\textbf{Additional Example is Powered by Desmos. Activities are available at SavvasRealize.com. ADDITIONAL EXAMPLE 5. ESSENTIAL QUESTION. Go back. SOLUTION.}
The table and graph show the average monthly rainfall for a city. Write the trigonometric function that models this data.
\begin{tabular}{ll}
Month & Rainfall(in) \\
1 & 6.4 \\
2 & 6 \\
3 & 6.3 \\
4 & 7.5 \\
5 & 9 \\
6 & 10.6 \\
7 & 11.4 \\
8 & 12 \\
9 & 11.5 \\
10 & 10.7 \\
11 & 9 \\
12 & 7.5 \\
\end{tabular}
% FIG 56 924 938 1128 1051
\placeholder{graph}{Orange scatterplot rainfall over two years, arrowed x,y axes. x labels0,2,...24, grid1; y labels0,2,...12, grid1, window0..24,0..13. First year table values at1..12; point(0,7.5), second-year repeated points13..24, maxima(8,12),(20,12), minima(2,6),(14,6).}
The midline is (\frac{12+6}2=9).
The amplitude is (\frac{12-6}2=3).
The period is (\frac\pi6).
The phase shift is5.
\answer{(y=3\sin(\frac\pi6(x-5))+9)}
\te{Printed period pi/6; likely12, with coefficient pi/6.}
Example5 Help students write a trigonometric function to model a real-world situation.
Q: How does the midline of the graph of the function compare with the average of the monthly rainfalls?
\answer{The midline value of9 is close to the value of the average of the monthly rainfalls,8.95.}
\end{te-addendum}
\begin{example}{1}
\textbf{Understand Phase Shift as a Horizontal Translation}
CONCEPTUAL UNDERSTANDING. How can you identify a horizontal translation from the rule of a trigonometric function?
A. How does changing the value of (c) affect the graph of (f(x)=2\sin(x-c))?
Sketch the graphs:
% FIG 56 854 317 1150 477
\placeholder{graph}{Two side-by-side graphs on arrowed x,y axes with origin O; ticks x pi/4, labeled-pi,pi; y ticks1/2 labeled-2,1,2, window about-5pi/4..5pi/4,-5/2..5/2. Left red f(x)=2sin x and blue f(x)=2sin(x+pi/2), blue maximum(0,2), red max(pi/2,2), blue left arrow bracket pi/2 units. Right red2sin x and green2sin(x-pi/6), green shifted rightpi/6 with green right bracket. All curves arrowheads. Callout: For each graph, the amplitude is2 and the period is2pi.}
\te{LOOK FOR RELATIONSHIPS. Just as with other functions you have studied, adding or subtracting a constant value from x as in f(x-c) shifts the graph of the function horizontally. Adding or subtracting a constant value from y as in f(x)-c shifts the graph of the function vertically.}
Changing (c) translates the graph left or right.
It does not change the amplitude or period.
A horizontal translation of a periodic function is called a phase shift.
\answer{A. Changing c translates the graph left or right; amplitude and period do not change.}
% Source page 57: 374 Example1 continued
B. How does changing the value of (c) affect the graph of (g(x)=\sin2(x-c))?
a. Sketch the graphs:
% FIG 57 792 266 1115 419
\placeholder{graph}{Two adjacent graphs with arrowed x,y axes,O; x ticks pi/4,labels-pi,pi, window-3pi/2..3pi/2; y ticks1/2,labels-2,1,2, window-5/2..5/2. Red g(x)=sin2x in both. Left blue g(x)=sin2(x+pi/2), reflection of red, with leftward pi/2 shift bracket. Right green g(x)=sin2(x-pi/6) with rightward pi/6 bracket. All curves arrowed, amplitude1 periodpi. Callouts: For these graphs, the amplitude is1. For these graphs, the period is pi.}
\te{COMMON ERROR. You may think that the period of (y=\sin2x) is twice the period of (y=\sin x), or (4\pi). In fact, the period of (y=\sin2x) is half the period of (y=\sin x), or (\pi).}
Changing (c) shifts the graph left or right as before.
The amplitude and period are the same for all three graphs.
\answer{B. Changing c shifts the graph left or right; amplitude1 and period pi are unchanged.}
\end{example}
\begin{tryit}{1}
Sketch the graph.
a. (h(x)=3\cos(x+\frac\pi4))
b. (h(x)=\cos(3x+\frac{3\pi}4))
\answer{a.,b. Graphs shown.}
% FIG 57 138 173 305 340
\placeholder{graph}{Try1a red3cos(x+pi/4) on square grid, arrowed x,y axes,O; x ticks pi/2,labels-pi,pi,window-2pi..2pi; y ticks1,labels-2,2,window-4..4. Maxima at-pi/4 and7pi/4,y3; minima-5pi/4,3pi/4,y-3; curve arrows.}
% FIG 57 349 173 543 341
\placeholder{graph}{Try1b red cos(3x+3pi/4) with amplitude1 period2pi/3; x,y arrowed axes,O; x ticks pi/2,labels-pi,pi,window-5pi/2..5pi/2; y ticks1/2,labels-1,1,window-2..2. Maxima x=-pi/4+2kpi/3 and minima pi/12+2kpi/3; curve end arrows.}
\end{tryit}
\begin{te-addendum}{2}{GeneralizeETP}
\textbf{Graph a Sine or Cosine Function. Use and Connect Mathematical Representations}
Q: How does the parameter (b) influence the period of the sine or cosine function?
\answer{The period is calculated by dividing (2\pi) by b.}
Q: How does the parameter (d) influence the graph of the sine or cosine function?
\answer{The parameter d shifts the graph of the function up or down d units. Adding d to the function adds d units to the output of the function.}
Examples are available at SavvasRealize.com.
\end{te-addendum}
\begin{example}{2}
\textbf{Graph a Sine or Cosine Function}
How can knowledge of transformations help you sketch the graph of (f(x)=2\sin\frac x2+3)?
\te{USE STRUCTURE. Make sure you understand which parameter influences each feature of the function (f(x)=a\sin b(x-c)+d). (|a|): amplitude; b: period; c: phase shift; (1/c): frequency; d: vertical shift.}
\te{Printed frequency1/c; likely b/(2pi).}
(f(x)=\sin x): Period(2\pi); Amplitude1; Phase shift0; Vertical shift0.
(f(x)=\sin\frac x2): Period(4\pi).
(f(x)=2\sin\frac x2): Period(4\pi); Amplitude2.
(f(x)=2\sin\frac x2+3): Period(4\pi); Amplitude2; Phase shift0; Vertical shift3.
% FIG 57 792 575 1115 775
\placeholder{graph}{Two stacked sine transformation diagrams with callout labels. Arrowed x,y axes,O; x ticks pi/2 labeled-pi,pi,window-3pi..3pi; y ticks1 labeled2,4,window about-3..6. Top red sin x amplitude1 period2pi and blue sin(x/2) amplitude1 period4pi. Bottom green2sin(x/2) min(-pi,-2),max(pi,2) and purple2sin(x/2)+3 min(-pi,1),max(pi,5); all curves have end arrows. Labels give each equation, period, amplitude, phase and vertical shifts as transcribed.}
\answer{(f(x)=2\sin\frac x2+3): period(4\pi), amplitude2, phase shift0, vertical shift3; graph shown.}
\end{example}
\begin{tryit}{2}
Sketch the graph of the function (g(x)=\frac23\cos(x-\frac\pi2)+1).
\answer{Graph shown.}
% FIG 57 125 649 280 805
\placeholder{graph}{Try2 red(2/3)cos(x-pi/2)+1 on square grid with arrowed x,y axes,O; x ticks pi/2 labels-pi,pi,window-2pi..2pi; y ticks1/2,labels1,2,window-1..3. Maxima(-3pi/2,5/3),(pi/2,5/3), minima(-pi/2,1/3),(3pi/2,1/3),yintercept1,end arrows.}
\end{tryit}
\begin{te-addendum}{2}{ElicitEvidence}
\textbf{Elicit and Use Evidence of Student Thinking}
Q: In the general form of the cosine function, (y=a\cos b(x-c)+d), what do the parameters (a), (b), (c), and (d) represent?
\answer{The a determines the amplitude, (2\pi) divided by b is the period, c is the phase shift, and d is the vertical shift.}
\end{te-addendum}
\begin{te-addendum}{2}{CommonError}
\textbf{Common Error. Try It!2}
Students may think that the phase shift of this function is (\pi/2) units to the left instead of to the right. Have students review the general form of the function (y=a\sin b(x-c)+d). Point out that (c) is subtracted from the value for (x), so in this case (c) is a positive number and the phase shift is to the right.
\end{te-addendum}
\begin{te-addendum}{2}{RtIExtend}
\textbf{Extend Student Thinking. USE WITH EXAMPLE2}
Have students graph a more complex sine or cosine function with (a), (b), (c), and (d\neq1).
\item Use transformations to help you graph (y=3\sin(3(x-2))+5).
Q: What are the amplitude, period, phase shift, and vertical shift?
\answer{The amplitude is3, the period is (2\pi/3), the phase shift is right2 units, and the vertical shift is up5 units.}
Q: Sketch the graph of this function.
\answer{Check students' graphs.}
Q: How can transformations help you sketch the graph of a function?
\answer{Understanding transformations eliminates the need to calculate the values of the function at key x-values in order to plot individual points of a transformed function.}
\end{te-addendum}
% Source page 58: 375 Understand & Apply
\begin{te-addendum}{3}{GeneralizeETP}
\textbf{Analyze a Sine or Cosine Function. Use and Connect Mathematical Representations}
PART A
Q: Why is it important to identify the key features of a sine or cosine function?
\answer{The key features can help you sketch the graph of a function without having to plot individual points.}
PART B
Q: Why is the amplitude positive when the sign of (a) is negative?
\answer{Amplitude is a measure of height. The amplitude is the absolute value of a because height cannot be negative.}
Q: What does the negative coefficient of the function mean?
\answer{The graph is a reflection of the parent function over the x-axis.}
Examples are available at SavvasRealize.com.
\end{te-addendum}
\begin{example}{3}
\textbf{Analyze a Sine or Cosine Function}
A. What are the key features of the graph of the function (f(x)=4\cos(3x)+1)?
(f(x)=4\cos3x+1)
(Amplitude=|4|=4)
\te{Callout: The amplitude is the absolute value of the coefficient of cosine.}
(Period=\frac{2\pi}3)
\te{Callout: To calculate the period, divide (2\pi) by the coefficient of x.}
Vertical shift=1
\te{Callout: The graph of (y=4\cos3x) has been shifted up1 unit.}
There is no phase shift. The graph has not been shifted horizontally from that of (y=\cos x). In terms of the general form, (y=a\cos[b(x-c)]+d), (c=0).
Sketching the graph can help make sense of these parameters.
The amplitude of (y=\cos x) is1. Notice that, for this function, the maximum value is (4+1), or5, and the minimum value is (-4+1), or-3.
% FIG 58 1034 389 1153 487
\placeholder{graph}{Red4cos3x+1, arrowed x,y axes,O, x ticks pi/6 labels-pi/2,pi/2; y ticks1 labels-2,2,4,window-pi..pi,-4..6. Maxima(-2pi/3,5),(0,5),(2pi/3,5), minima(-pi/3,-3),(pi/3,-3),curved end arrows.}
\te{USE APPROPRIATE TOOLS. Why is it important to identify the key features of a function, even if you intend to graph it on a calculator?}
\answer{A. Amplitude4,period(2\pi/3),vertical shift1,phase shift0,maximum5,minimum-3.}
B. What are the key features of the graph of the function (f(x)=-3\sin(x+\frac\pi2)-2)?
(f(x)=-3\sin(x+\frac\pi2)-2)
(Amplitude=|-3|=3)
(Period=2\pi)
Vertical shift=2 units down
Phase shift=(\pi/2) units left
\te{Callout: Since the coefficient of the sine function is negative, this graph is a reflection over the x-axis of the parent function.}
Sketch the graph to check.
Notice that the combined effects of the reflection and the phase shift produce a graph that could be interpreted as a sine graph without a reflection and with a phase shift (\pi/2) units right.
% FIG 58 969 596 1153 695
\placeholder{graph}{Red-3sin(x+pi/2)-2, arrowed x,y axes,O, x labels pi/2,pi,3pi/2,2pi,gridpi/8; y labels-6,-4,-2,2,grid1,window-pi/4..17pi/8,-7..3. Minima(0,-5),(2pi,-5), maximum(pi,1),end arrows.}
\answer{B. Amplitude3,period(2\pi),vertical shift2 down,phase shift(pi/2) left.}
\end{example}
\begin{tryit}{3}
Identify the amplitude, period, phase shift, vertical shift, and the maximum and minimum values of each function.
a. (g(x)=\frac12\sin(x-\frac\pi3)-4)
\answer{a. amplitude:(1/2); period:(2\pi); phase shift:(pi/3) units to the right; vertical shift:4 units down; minimum value:-4.5; maximum value:-3.5}
b. (g(x)=2\cos(2x+\frac\pi4)+2)
\answer{b. amplitude:2; period:(\pi); phase shift:(pi/8) to the left; vertical shift:2 units up; minimum value:0; maximum value:4.}
\end{tryit}
\begin{te-addendum}{3}{HabitsOfMind}
\textbf{HABITS OF MIND. Use with EXAMPLES1,2,\&3}
Look for Relationships For other types of functions, (y=f(x-a)) was always a horizontal translation by (a) units of the parent function (y=f(x)). Does that rule hold for trigonometric functions? Explain.
\answer{Yes, when you graph (y=\frac23\cos(x-\frac\pi2)+1), it is shifted right (\pi/2) units.}
\end{te-addendum}
\begin{te-addendum}{3}{RtISupport}
\textbf{Support Struggling Students. USE WITH EXAMPLE3}
Some students may struggle with analyzing sine functions. Have students use this problem for extra practice.
\item What are the key features of the graph of (y=\frac12\sin(x-4)-3)?
Q: What is the amplitude of this function? Explain.
\answer{(1/2); the amplitude is the absolute value of the coefficient of the sine function.}
Q: What is the phase shift? Explain.
\answer{Right4 units; the phase shift is the value subtracted from x.}
Q: What is the vertical shift?
\answer{down3 units}
Q: Sketch the graph of the function.
% FIG 58 695 1073 1002 1270
\placeholder{graph}{Red(1/2)sin(x-4)-3 on grid, x,y arrowed axes,O; x ticks pi/4 labels-3pi/2,-pi,-pi/2,pi/2,pi,3pi/2; y ticks1 labels-4,-2,2,4,window-2pi..2pi,-5..5. Red maxima near(-.71,-2.5),(5.57,-2.5),minima(-3.85,-3.5),(2.43,-3.5),curve end arrows.}
\answer{Graph shown.}
\end{te-addendum}
% Source page 59: 376 Understand & Apply
\begin{te-addendum}{4}{GeneralizeETP}
\textbf{Write the Equation of a Translation. Use and Connect Mathematical Representations}
Q: What do you notice about the graph of the stamp's height?
\answer{It is similar to the graph of (y=\sin x) because its period is (2\pi) and its amplitude is1.}
Q: How could you determine that the period is (2\pi) even though the x-axis is given in seconds?
\answer{The graph crosses the x-axis at0 and at slightly more than6. If you divide6 by pi it is approximately2.}
Examples are available at SavvasRealize.com.
\end{te-addendum}
\begin{example}{4}
\textbf{Write the Equation of a Translation}
APPLICATION. An assembly-line machine stamps a company's logo on each product coming down the line. The graph models the stamp's height from the machine's axle, (y), at time (x) seconds since the machine turns on. What is an equation that models this motion? Interpret the graph's y-intercept in terms of the context.
% FIG 59 923 240 1080 300
\placeholder{graph}{Red sinusoid on grid, arrowed x,y axes,O; x labels-2,2,4,6,8,10,12,grid1; y labels-1,1,grid1/2;window-3..13,-3/2..3/2. Red filled point labeled(3,0),upward crossing; troughs around1.5,7.8 at-1,peaks4.6,10.9 at1;curve startsnearorigin movingdown and arrowatright.}
% FIG 59 922 282 1113 378
\placeholder{illustration}{Two gray robotic stamping arms beside a brown factory conveyor carrying yellow cardboard boxes with blue labels, in perspective.}
\te{STUDY TIP. Because sine and cosine are simply phase shifts of each other, you can build the equation using either function.}
The sine function has a general form
(f(x)=a\sin b(x-c)+d).
The amplitude is1. You may be tempted to define this graph as reflected about the x-axis, but it is easier to work with just a phase shift instead of a phase shift and a reflection. (a=1).
This graph contains the point (3,0), where the corresponding point on (y=\sin x) would be (0,0), this is a horizontal shift of3 to the right, so (c=3).
The period is (2\pi), so (b=1).
There is no vertical shift---the midline is the y-axis, just like the midline of (y=\sin x). So (d=0).
\te{Printed midline y-axis; likely x-axis.}
So an equation for the graph is (g(x)=\sin(x-3)).
At (x=0), when the equipment is turned on, the equipment is near the midline of its cycle, moving down.
\answer{(g(x)=\sin(x-3)); at x=0 the equipment is near the midline of its cycle,moving down.}
\end{example}
\begin{tryit}{4}
Write an equation that models the function represented by the graph.
% FIG 59 852 614 971 665
\placeholder{graph}{Blue cosine minus2 on grid, arrowed x,y axes,O; x ticks pi/2 labeled-2pi,-pi,pi,2pi; y ticks1 labeled-2,window-3pi..3pi,-4..1. Blue filled maximum(0,-1) labeled,other maxima(+-2pi,-1),minima(+-pi,-3);end arrows.}
\answer{(y=\cos(x)-2)}
\end{tryit}
\begin{te-addendum}{4}{ElicitEvidence}
\textbf{Elicit and Use Evidence of Student Thinking}
Q: What key features can you identify by looking at the graph?
\answer{It has an amplitude of1; it is shifted down2 units; it is a cosine function or a sine function that has been phase shifted; it has a period of (2\pi).}
\end{te-addendum}
\begin{te-addendum}{4}{ELLAddendum}
\textbf{English Language Learners. Use with EXAMPLE4}
READING. BEGINNING. Ask students to read the Study Tip for the example.
Q: How is the word phase used in the sentence?
\answer{It is used as an adjective to define the type of shift between sine and cosine.}
Q: What is another word that could be used instead of phase to describe the kind of shift?
\answer{horizontal}
Q: What is another word that could be used instead of shift?
\answer{movement}
WRITING. DEVELOPING. Have students write the definitions of corresponding and point in their journals. Then have students answer the following questions in their journals.
Q: If the points (3,0) and (0,0) are corresponding, then why do they not have the same values for (x) and (y)?
\answer{They are on the same place on the curve, but that spot on the curve occurs at different locations on the x-axis.}
Q: What is another way that the word corresponding can be used?
\answer{as a way of communicating, such as by text or email.}
LISTENING. EXPANDING. Read the first two sentences of the example aloud. Discuss the different usages of the word stamp, as a verb and a possessive noun. Have students listen to the following sentence pairs and determine which form of the italicized word is used for each sentence.
Q: (1) The circle's area is larger than the triangle's area. (2) Brian circles the correct answer.
\answer{possessive noun; verb}
Q: (1) Daquan graphs the equation on the board. (2) The year of birth is represented by the graph's x-axis.
\answer{verb; possessive noun}
\end{te-addendum}
% Source page 60: 377 Understand & Apply
\begin{te-addendum}{5}{GeneralizeETP}
\textbf{Find a Trigonometric Model. Use and Connect Mathematical Representations}
Q: What type of function could be used to model this graph?
\answer{Either a sine or cosine function could be used.}
Q: What can you identify by looking at the graph?
\answer{The graph has a vertical shift, a phase shift, and an increased amplitude.}
Q: How do you determine the phase shift of the function?
\answer{Look for the point where the graph crosses the midline to determine the phase shift.}
Q: How do you determine the vertical shift of the function?
\answer{The vertical shift is found by calculating the average of the maximum and the minimum.}
Examples are available at SavvasRealize.com.
\end{te-addendum}
\begin{example}{5}
\textbf{Find a Trigonometric Model}
APPLICATION. The table shows the average high temperature by month for Washington, DC. How can these temperatures be modeled with a trigonometric graph? How does the midline function value compare with the average of the12 temperatures?
\begin{tabular}{lllllllllllll}
Month & Jan. & Feb. & Mar. & April & May & June & July & Aug. & Sept. & Oct. & Nov. & Dec. \\
High(°F) & 43 & 47 & 56 & 67 & 75 & 84 & 88 & 87 & 80 & 68 & 58 & 47 \\
\end{tabular}
Formulate. Assign numbers to represent the months of the year, and create a scatter plot to show the pattern of temperatures in the table repeating over a span of two years.
% FIG 60 838 368 1085 462
\placeholder{graph}{Washington D.C. Average High Temperature(°F), blue scatterplot x,y arrowed axes; x labels0,4,8,12,16,20,24,grid1; y labels30,50,70,90,grid10,window0..24,30..100. Points at integer months1..12 with y43,47,56,67,75,84,88,87,80,68,58,47, repeated at13..24.}
The temperatures can be modeled by a function in the form (y=a\sin b(x-c)+d).
Compute. Average the maximum and minimum temperatures to find the midline (vertical shift). (\frac{88+43}2=65.5), so the vertical shift is (d=65.5).
Find the amplitude.
(amplitude=\frac{\max-\min}2=\frac{88-43}2=22.5)
so the amplitude (a=22.5)
The period is12 months, so (b=\frac{2\pi}{12}=\frac\pi6).
Estimate the phase shift: (c\approx4)
\te{Callout: Determine the phase shift by looking for the point where the graph crosses the midline. In the parent function, there is a zero exactly at the origin.}
(y=22.5\sin[\frac\pi6(x-4)]+65.5)
Interpret. The average of the12 temperatures,66.7°F, is close to the midline value,65.5°F.
\answer{(y=22.5\sin[\frac\pi6(x-4)]+65.5); average66.7°F close to midline65.5°F.}
\end{example}
\begin{tryit}{5}
Write a trigonometric function to model the average high temperatures for Philadelphia, Pennsylvania. How does the midline value compare with the average of the12 temperatures?
\begin{tabular}{lllllllllllll}
Month & Jan. & Feb. & Mar. & April & May & June & July & Aug. & Sept. & Oct. & Nov. & Dec. \\
High(°F) & 40 & 44 & 53 & 64 & 74 & 83 & 87 & 85 & 78 & 67 & 56 & 45 \\
\end{tabular}
\answer{(y=23.5\sin(\frac\pi6(x-4))+63.5); The average of the12 temperatures,64.7°F,is close to the midline value,63.5°F.}
\end{tryit}
\begin{te-addendum}{5}{ElicitEvidence}
\textbf{Elicit and Use Evidence of Student Thinking}
Q: What is the first step to help you write the equation for this data set?
\answer{Create a scatterplot of the data points and compare it to the graph of the sine or cosine function.}
\end{te-addendum}
\begin{te-addendum}{5}{HabitsOfMind}
\textbf{HABITS OF MIND. Use with EXAMPLES4\&5}
Model with Mathematics How do you know whether you should use a sine or cosine function to model a periodic graph?
\answer{You could use either function, but if there is a zero at the y-axis, it might make more sense to use sine, whereas if there is a relative maximum at the y-axis, you might prefer to use cosine.}
\end{te-addendum}
% Source page 61: 378 Concept Summary
\begin{te-addendum}{ConceptSummary}{PurposefulQuestions}
\textbf{CONCEPT SUMMARY. Transformations of Trigonometric Functions}
Q: How does knowing the general form of a function help you understand and graph transformations of trigonometric functions?
\answer{If you know the general form of the function,(y=a\sin b(x-c)+d),and what the parameters a,b,c,and d represent,then you can identify key features of the graph of the function,such as amplitude,period,phase shift,and vertical shift.}
Concept Summary, Do You Understand, and Do You Know How are available at SavvasRealize.com. Presentation. Practice.
\end{te-addendum}
\begin{te-addendum}{ConceptSummary}{CommonError}
\textbf{Do You UNDERSTAND? Do You KNOW HOW? Common Error. Exercise5}
Some students may mistakenly think that the phase shift is (\pi/6) units to the left because the expression in the input is (x-\frac\pi6). Have students look at the general form of the function again to notice that the value for the phase shift is subtracted from (x), so the value of the phase shift is positive in this case, which means a shift to the right.
\end{te-addendum}
\begin{concept-summary}
\textbf{Transformations of Trigonometric Functions}
GRAPH
% FIG 61 697 260 866 362
\placeholder{graph}{Red y=sin x and blue y=2sin[(1/2)(x-pi/2)]+1,arrowed x,y axes,O; x ticks pi/2,labels-2pi,-pi,pi,2pi; y ticks1,labels-2,2,window-3pi..3pi,-4..4. Blue maxima(-5pi/2,3),(3pi/2,3),minimum(-pi/2,-1), red amplitude1 period2pi;curves end arrows,equation labels colored.}
WORDS AND SYMBOLS
(y=a\sin b(x-c)+d)
(y=2\sin[\frac12(x-\frac\pi2)]+1)
(Amplitude=|a|=|2|=2)
Stretches the graph by a factor of2
(Period=\frac{2\pi}b=\frac{2\pi}{1/2}=4\pi)
Cycle repeats every (4\pi) units.
(Phase\ Shift=c=\frac\pi2)
Graph is shifted (\pi/2) units to the right.
(Vertical\ Shift=d=1)
Graph is shifted up one unit.
\textbf{Do You UNDERSTAND?}
1. ESSENTIAL QUESTION How can you find and use translations of graphs of trigonometric functions?
\answer{You can analyze the graphs of sine and cosine functions to find changes in amplitude,period,phase shift,and vertical shift.}
2. Vocabulary What is a phase shift?
\answer{A phase shift is a horizontal translation of a periodic function.}
3. Error Analysis Felipe said the function (y=\frac12\cos[3(x+\frac\pi4)]-3) has a phase shift (\pi/4) units to the right and a vertical shift3 units down. Describe and correct the error Felipe made.
\answer{The value of c in the function is negative,which means the phase shift is to the left,not to the right. The phase shift is left(pi/4) units.}
4. Use Structure Write a sine function that has an amplitude of (1/6),a period of (8\pi/3),a phase shift of (2\pi) units to the right,and a vertical shift of5 units up.
\answer{(y=\frac16\sin[\frac34(x-2\pi)]+5)}
\textbf{Do You KNOW HOW?}
Identify the amplitude,period,phase shift,and vertical shift of the function.
5. (y=4\sin(x-\frac\pi6)+2)
\answer{amplitude:4;period:(2\pi);phase shift:right(pi/6) units;vertical shift:up2 units}
6. (y=\frac13\cos[2(x+\frac\pi2)]-1)
\answer{amplitude:(1/3);period:(\pi);phase shift:left(pi/2) units;vertical shift:down1 unit}
7. Write an equation for the function represented by the graph using the cosine function.
% FIG 61 919 592 1039 681
\placeholder{graph}{Red sinusoid,arrowed x,y axes,O; x ticks pi/2 labels-2pi,-pi,pi,2pi; y ticks1/2 labels-1,-2,window-3pi..3pi,-3..1. Red labeled points(-pi/2,1),(0,-1),(pi/2,-3); maxima-5pi/2,-pi/2,3pi/2,y1;troughs-3pi/2,pi/2,5pi/2,y-3;arrows.}
\answer{(y=2\cos(x+\frac\pi2)-1)}
8. Sketch a graph of the function (y=\sin[2(x+\frac\pi2)]+1).
\answer{Graph shown.}
% FIG 61 126 994 282 1187
\placeholder{graph}{Red-sin2x+1,arrowed x,y axes,O; x ticks pi/2 labels-pi,pi; y ticks1/2 labels-2,-1,1,2;window-2pi..2pi,-5/2..5/2. Maxima(-pi/4+kpi,2), minima(pi/4+kpi,0),yintercept1,end arrows.}
\end{concept-summary}
% Source page 62: 379 Practice & Problem Solving
\begin{te-addendum}{Practice}{GeneralizeETP}
\textbf{PRACTICE \& PROBLEM SOLVING}
Practice \& Problem Solving and video tutorials are available at SavvasRealize.com.
Lesson Practice You may opt to have students complete the automatically scored Practice and Problem Solving items online powered by MathXL for School.
Choose from: Lesson Practice; Adaptive Practice
You may also take advantage of the bank of exercises for assigning additional practice.
\textbf{Assignment Guide}
\begin{tabular}{ll}
On-Level & Advanced \\
9--20,22--28 & 9--15,16--28 \\
\end{tabular}
\textbf{Engage Through Student Choice}
Promote student agency by allowing students to choose practice items. You may structure this choice in many ways. For example:
Assign each section a point value. Students choose at least one item from each section and items chosen should have a minimum of20 total points.
Understand, Apply:2 points each
Practice:1 point each
Assessment Practice:1 point each
Performance Task:3 points each
Scan the QR code to access a video tutorial.
\end{te-addendum}
\begin{item-analysis}
\begin{tabular}{lll}
\toprule
Example & Items & DOK \\
\midrule
1 & 9,10,14,16,17 & 2 \\
2 & 11,18,19 & 2 \\
3 & 12,15,20,21,26,27 & 2 \\
3 & 28 & 3 \\
4 & 13,22 & 2 \\
5 & 23--25 & 2 \\
\bottomrule
\end{tabular}
\end{item-analysis}
\begin{practice}{9}{2}
UNDERSTAND. Use Structure Write a sine function and a cosine function for the graph.
% FIG 62 551 318 780 447
\placeholder{graph}{Red sinusoid on grid,arrowed x,y axes,O; x ticks pi/4 labels-2pi,-pi,pi,2pi; y ticks1/4 labels-1,1;window-9pi/4..9pi/4,-5/4..5/4. Marked points(-pi/6,1),(pi/3,0),(5pi/6,-1), red labels;period2pi,end arrows.}
\answer{(y=\sin(x+\frac{2\pi}3)) and (y=\cos(x+\frac\pi6))}
\end{practice}
\begin{practice}{10}{2}
UNDERSTAND. Error Analysis Describe and correct the error a student made in finding the phase shift of the given function.
\te{Student work marked with redX:}
(y=\cos(3x+\frac\pi2))
phase shift=(\pi/2) units to the left
\answer{The student did not first factor out a3. The phase shift is left(pi/6) units.}
\end{practice}
\begin{practice}{11}{2}
UNDERSTAND. Generalize Describe the phase shift and vertical shift of a function in the form (y=a\sin[b(x-c)]+d).
\answer{When(c>0),the phase shift is right c units. When(c<0),the phase shift is left|c| units. When(d>0),the vertical shift is up d units. When(d<0),the vertical shift is down d units.}
\te{Printed down d units for negative d; likely down|d| units.}
\end{practice}
\begin{practice}{12}{2}
UNDERSTAND. Higher Order Thinking How are the domain and range of the function (y=\frac14\cos[3(x-\frac{2\pi}3)]+2) related to the domain and range of the parent function (y=\cos x)? Explain your reasoning.
\answer{The graph of(y=\frac14\cos[3(x-\frac{2\pi}3)]+2) is compressed by a factor of(1/4),shifted right(2\pi/3) units,and shifted up2 units from the graph of(y=\cos x). The horizontal shift does not affect the domain and range of the function. The compression factor divides the range interval by one-fourth and the vertical shift up2 units adds2 units to the compressed range.}
\te{Printed divides by one-fourth; likely multiplies by one-fourth.}
\end{practice}
\begin{practice}{13}{2}
UNDERSTAND. Generalize Write an equation for the midline of the function (y=a\cos[b(x-c)]+d).
\answer{(y=d)}
\end{practice}
\begin{practice}{14}{2}
UNDERSTAND. Reason How are the zeros of the function (y=\sin(x+\frac\pi3)) related to the zeros of the parent function (y=\sin x)?
\answer{The zeros of(y=\sin(x+\frac\pi3)) are shifted left(pi/3) units of the zeros of(y=\sin x).}
\end{practice}
\begin{practice}{15}{2}
UNDERSTAND. Mathematical Connections In the equation (y=a\sin[b(x-c)]+d),which of the parameters (a),(b),(c),and (d) can have an effect on the y-intercept of the graph? Explain.
\answer{Answers may vary. Sample: The y-intercept is always affected by the value of d. When(c\neq0),the phase shift means that a,b,and c can also affect the y-intercept by changing the amplitude and period.}
\end{practice}
\begin{practice}{16}{2}
PRACTICE. Sketch the graph of the function. SEE EXAMPLE1
(y=\cos(x-\frac\pi4))
\answer{Graph shown.}
% FIG 62 853 1176 1010 1332
\placeholder{graph}{Red cos(x-pi/4) on square grid,arrowed x,y axes,O; x ticks pi/2 labeled-pi,pi; y ticks1/2 labeled-1,1;window-2pi..2pi,-2..2. Max(pi/4,1),min(-3pi/4,-1),(5pi/4,-1),end arrows.}
\end{practice}
\begin{practice}{17}{2}
PRACTICE. Sketch the graph of the function. SEE EXAMPLE1
(y=2\sin(x+\frac{3\pi}4))
\answer{Graph shown.}
% FIG 63 125 167 281 325
\placeholder{graph}{Red2sin(x+3pi/4),x,y arrowed axes,O; x gridpi/2 labeled-pi,pi; y grid1/2 labeled-1,1,window-2pi..2pi,-2..2. Max(-pi/4,2),(7pi/4,2),min(-5pi/4,-2),(3pi/4,-2),end arrows.}
\end{practice}
\begin{practice}{18}{2}
PRACTICE. Sketch the graph of the function. SEE EXAMPLE2
(y=\frac13\cos(x+\frac\pi2)-2)
\answer{Graph shown.}
% FIG 63 125 342 280 498
\placeholder{graph}{Red(1/3)cos(x+pi/2)-2,x,y arrowed axes,O; x gridpi/2 labeled-pi,pi; y grid1/2 labeled-1,-2,window-2pi..2pi,-3..1. Max(-pi/2,-5/3),(3pi/2,-5/3),min(-3pi/2,-7/3),(pi/2,-7/3),end arrows.}
\end{practice}
\begin{practice}{19}{2}
PRACTICE. Sketch the graph of the function. SEE EXAMPLE2
(y=3\sin(x-\frac\pi6)+1)
\answer{Graph shown.}
% FIG 63 125 517 282 674
\placeholder{graph}{Red3sin(x-pi/6)+1,arrowed x,y axes,O; x gridpi labeled-2pi,2pi; y grid1 labeled-2,2,4,window-4pi..4pi,-3..5. Max2pi/3+2kpi,y4;min5pi/3+2kpi,y-2;fourcycles,end arrows.}
\end{practice}
\begin{practice}{20}{2}
PRACTICE. Identify the amplitude,period,phase shift,vertical shift,and the maximum and minimum values of the function. SEE EXAMPLE3
(y=\frac23\sin(x+\frac\pi3)+3)
\answer{amplitude:(2/3);period:(2\pi);phase shift:left(pi/3) units;vertical shift:up3 units;maximum:(3\frac23);minimum:(2\frac13)}
\end{practice}
\begin{practice}{21}{2}
PRACTICE. Identify the amplitude,period,phase shift,vertical shift,and the maximum and minimum values of the function. SEE EXAMPLE3
(y=\frac12\cos[2(x-\frac\pi4)]-1)
\answer{amplitude:(1/2);period:(\pi);phase shift:right(pi/4) units;vertical shift:down1 unit;maximum:(-1/2);minimum:(-1\frac12)}
\end{practice}
\begin{practice}{22}{2}
PRACTICE. Write an equation for the function represented by the graph using the sine function. SEE EXAMPLE4
% FIG 62 839 493 1118 571
\placeholder{graph}{Red sinusoid on grid,arrowed x,y axes,O; x gridpi/8 labels-pi,-pi/2,pi/2,pi,3pi/2,2pi; y grid1/2 labeled2;window-9pi/8..17pi/8,-1..5/2. Red marked points(0,3/2),(pi/2,1),(pi,1/2),period2pi,midline1,amplitude1/2,end arrows.}
\answer{(y=\frac12\sin(x-\frac{3\pi}2)+1)}
\end{practice}
\begin{practice}{23}{2}
PRACTICE. The table shows the brightness of the moon at the end of eight consecutive weeks. How can you model this with a trigonometric function? How does the midline of the function compare with the average of the8 visibility levels? SEE EXAMPLE5
\begin{tabular}{ll}
Week & Percent Visible \\
1 & 50\% \\
2 & 0\% \\
3 & 48\% \\
4 & 100\% \\
5 & 67\% \\
6 & 5\% \\
7 & 34\% \\
8 & 95\% \\
\end{tabular}
% FIG 62 1006 709 1091 1001
\placeholder{illustration}{Eight stacked moon-phase images on black starry backgrounds aligned with weeks1..8: left half illuminated;dark disk;right half illuminated;full moon;mostly left illuminated;thin left crescent;right crescent aboutthird;nearlyfull.}
\answer{(y=50\cos(\frac{7\pi}{15}x)+50); The average amount of moon visible during the8 weeks,about49.9\%,is nearly the same as the midline value of the function,50.}
\end{practice}
% Source page 63: 380 Practice & Problem Solving
\begin{practice}{24}{2}
APPLY. Model With Mathematics Alternating current is the flow of charge that periodically changes direction. Alternating current is used to deliver power. The function (V(t)=E\cos(wt+\frac\pi2)) gives the voltage in amps for (t) seconds.
\te{Printed voltage in amps; likely volts.}
a. Erasto wants to find the voltage when (E=40) volts and (w=188) radians per second. Write a function to represent this situation.
\answer{a. (V(t)=40\cos(188t+\frac\pi2))}
b. Rewrite the function so that the coefficient of (t) is1.
\answer{b. (V(t)=40\cos[188(t+\frac\pi{376})])}
c. What is the amplitude of the function?
\answer{c.40}
d. What is the period of the function?
\answer{d.(\pi/94)}
e. What is the phase shift of the function?
\answer{e.shift left(pi/376) units}
f. Graph the function.
\answer{f.Graph shown.}
% FIG 63 143 1041 336 1237
\placeholder{graph}{Red40cos(188t+pi/2) on grid,arrowed axes horizontal x,vertical y,O; x labels-pi/48,pi/48,gridpi/144; y labels-40,-20,20,40,grid10;window-5pi/144..5pi/144,-50..50. Amplitude40,periodpi/94,descending throughorigin,sixcyclesapprox,end arrows.}
\end{practice}
\begin{practice}{25}{2}
APPLY. Make Sense and Persevere The table shows the average amount of rainfall in inches by month for Junction City, California.
\begin{tabular}{ll}
Month & Rainfall(in.) \\
January & 6.46 \\
February & 5.83 \\
March & 4.84 \\
April & 2.52 \\
May & 1.81 \\
June & 0.79 \\
July & 0.29 \\
August & 0.16 \\
September & 0.59 \\
October & 2.28 \\
November & 5.39 \\
December & 7.87 \\
\end{tabular}
% FIG 63 625 600 780 780
\placeholder{illustration}{Cars in several lanes of a highway in rain,gray mountains and bridge background;yellow,orange,red,blue,white cars and rain streaks.}
a. How can these rainfall amounts be modeled with a trigonometric graph?
\answer{a.Answers may vary.Sample:(y=3.855\sin[\frac\pi6(x-10.5)]+4.015)}
b. How does the midline function value compare with the average of the12 rainfall amounts?
\answer{b.The average of the12 rainfall amounts,about3.236 inches,is close to the midline value,4.015 inches.}
c. Graph the function.
\answer{c.Graph shown.}
% FIG 63 498 1137 711 1369
\placeholder{graph}{Red rainfall model,arrowed x,y axes,O; x labels4,8,12,16,grid2; y labels2,4,6,8,grid1;window-2..20,-2..10. Peaks near(1.5,7.87),(13.5,7.87),trough(7.5,.16),(19.5,.16),end arrows.}
\end{practice}
\begin{practice}{26}{2}
ASSESSMENT PRACTICE. Determine if each statement about the function (y=\frac34\cos[3(x+\frac\pi6)]-5) is true. Write yes or no.
A. The amplitude is (3/4).
B. The period is3.
C. The phase shift is (\pi/6) units to the right.
D. The vertical shift is5 units down.
\answer{A.yes;B.no;C.no;D.yes}
\end{practice}
\begin{practice}{27}{2}
ASSESSMENT PRACTICE. SAT/ACT Kadia graphed the function (y=2\sin(x+\frac\pi2)-1) but forgot to label the x-axis. What is the value of (d) on the x-axis?
% FIG 63 826 478 1030 646
\placeholder{graph}{Red2cos x-1 on grid,arrowed x,y axes,O; x letters a,b,c,d,e,f,g at successive pi/2 positions except origin: a=-pi,b=-pi/2,c=pi/2,d=pi,e=3pi/2,f=2pi,g=5pi/2. x gridpi/4;y labels-3,-2,-1,1,2,grid1/2;window-5pi/4..11pi/4,-4..5/2. Max(0,1),(2pi,1),min(-pi,-3),(pi,-3),end arrows.}
A. (\pi/2)
B. (\pi)
C. (3\pi/2)
D. (2\pi)
\answer{B}
\end{practice}
\begin{practice}{28}{3}
ASSESSMENT PRACTICE. Performance Task Marco is investigating phase shifts of the parent sine function,(y=\sin x). He wants to map the sine function onto itself.
Part A Write an equation of a function that has an identical graph but includes a phase shift.
\answer{Part A Sample:(y=\sin(x+2\pi))}
Part B Write an equation that will map the parent sine function onto itself by shifting the parent function to the right.
\answer{Part B Sample:(y=\sin(x-2\pi))}
Part C What do the equations in part(a) and part(b) tell you about the period of the sine function?
\answer{Part C The period of the sine function is(2\pi).}
Part D How many equations can you write to map the parent sine function onto itself? Explain.
\answer{Part D infinitely many;The sine function has a period of(2\pi),so the graph repeats on every interval of(2\pi). There are an infinite number of intervals of(2\pi),so there are an infinite number of equations that will map the sine function onto itself.}
\end{practice}
% Source page 64: 380A Assess & Differentiate
\begin{te-addendum}{Assess}{RtISupport}
\textbf{LESSON QUIZ}
Use the Lesson Quiz to assess students' understanding of the mathematics in the lesson.
Students can take the Lesson Quiz online or you can download a printable copy from SavvasRealize.com. The Lesson Quiz is also available in the Assessment Sourcebook.
Lesson Quiz is available at SavvasRealize.com. Assessment.
\textbf{Item Analysis}
\begin{tabular}{lll}
Item & DOK & Skills Review and Practice \\
1 & 1 & F89 \\
2 & 2 & F89 \\
3 & 2 & F91 \\
4 & 1 & F89 \\
5 & 2 & F89 \\
\end{tabular}
Use the student scores on the Lesson Quiz to prescribe differentiated assignments.
If students take the Lesson Quiz online,it will be automatically scored and appropriate differentiated practice will be assigned based on student performance.
\begin{tabular}{lll}
Intervention & 0--3 points & Reteach to Build Understanding;Mathematical Literacy and Vocabulary;Lesson Virtual Nerd videos \\
On-Level & 4 points & Enrichment \\
Advanced & 5 points & Enrichment \\
\end{tabular}
\end{te-addendum}
\begin{te-addendum}{Assess}{GeneralizeETP}
\textbf{7-4 Lesson Quiz: Translating Trigonometric Functions}
ASSESSMENT RESOURCES. Name \underline{\hspace{2cm}}. enVision Algebra2. SavvasRealize.com.
1. Select all the correct statements about the graphs of (y=\frac12\sin x) and (y=\frac12\sin(x-\frac\pi2)).
A. The graph of (y=\frac12\sin(x-\frac\pi2)) is shifted (\pi/2) units left from the graph of (y=\frac12\sin x).
B. The graph of (y=\frac12\sin(x-\frac\pi2)) is shifted (\pi/2) units right from the graph of (y=\frac12\sin x).
C. The period of (y=\frac12\sin(x-\frac\pi2)) is (\pi/2),and the period of (y=\frac12\sin x) is (2\pi).
D. The amplitude of (y=\frac12\sin(x-\frac\pi2)) is half the amplitude of (y=\frac12\sin x).
E. The amplitude of (y=\frac12\sin(x-\frac\pi2)) is the same as the amplitude of (y=\frac12\sin x).
\answer{1.B,E}
2. Function (f) is a cosine function with period (8\pi),amplitude3,and a local maximum at (f(0)=-3). What is the equation of the midline of the graph of (f)?
(y=\underline{\hspace{1cm}})
\answer{2.(y=-6)}
3. Which trigonometric function best models the average low temperatures shown in the scatterplot?
% FIG 64 980 555 1111 660
\placeholder{graph}{Average Low Temperature(°F),black scatterplot,arrowed x,y axes;horizontal Months; x labels0,4,8,12,16,20,24,grid2; y labels0,10,20,30,grid5,window0..24,0..35. Points integer0..24 follow12.5sin[pi(x-4)/6]+17.5 approximately:6.7,5,6.7,11.3,17.5,23.8,28.3,30,28.3,23.8,17.5,11.3,6.7 then repeat;maxmonths7,19,min1,13.}
A. (y=12.5\sin[\frac\pi6(x+4)]+17.5)
B. (y=17.5\sin(\frac\pi6 x)+12.5)
C. (y=12.5\sin(\frac\pi6 x)+17.5)
D. (y=12.5\sin[\frac\pi6(x-4)]+17.5)
\answer{3.D}
4. What is the phase shift of the graph of (y=-2\sin(\frac32 x-\frac\pi2)-4)?
A. (\pi/2) units right
B. (\pi/3) units right
C.4 units left
D. (3\pi/4) units right
\answer{4.B}
5. Select all the statements that describe the key features of the graph of (y=2\cos[\frac12(x+\frac\pi4)]+5).
A.amplitude=5
B.period=(4\pi)
C.phase shift is (\pi/4) units left
D.phase shift is5 units left
E.vertical shift is (\pi/4) units up
F.vertical shift is5 units up
\answer{5.B,C,F}
enVision Algebra2. Assessment Sourcebook. Copyright©Savvas Learning Company LLC.All Rights Reserved.
\end{te-addendum}
% Source page 65: 380B Differentiated Resources
\begin{te-addendum}{Assess}{GeneralizeETP}
\textbf{DIFFERENTIATED RESOURCES}
I=Intervention. O=On-Level. A=Advanced.
Reteach to Build Understanding I. Provides scaffolded reteaching for the key lesson concepts. MathXL for School.
Additional Practice I O. Provides extra practice for each lesson. MathXL for School.
Enrichment O A. Presents engaging problems and activities that extend the lesson concepts. MathXL for School.
Mathematical Literacy and Vocabulary I O. Helps students develop and reinforce understanding of key terms and concepts.
\end{te-addendum}
\begin{te-addendum}{Assess}{RtISupport}
\textbf{7-4 Reteach to Build Understanding. Translating Trigonometric Functions}
Name \underline{\hspace{2cm}}. enVision Algebra2. SavvasRealize.com.
1. Using the first two equations as a guide,complete the table. Remember that the period of the sine and cosine functions is (2\pi).
\begin{tabular}{lllll}
 & (y=a\sin b(x-c)+d) & (y=a\cos b(x-c)+d) & (y=\sin x) & (y=3\sin2(x-1)+6) \\
Period & (2\pi/b) & (2\pi/b) & \answer{(2\pi)} & (\pi) \\
Amplitude & (|a|) & (|a|) & 1 & \answer{3} \\
Phase Shift & c & c & \answer{0} & 1 \\
Vertical Shift & d & d & 0 & \answer{6} \\
\end{tabular}
2. Emaan stated that the graph shown of (y=2\sin(x-2)+3) is (y=2\sin(x)) translated left2 and up3. Is Emaan correct? Use the graph to help you describe the translation of the graph and correct Emaan's error.
\answer{Emaan is not correct. The graph was translated right2 and up3.}
% FIG 65 326 500 389 569
\placeholder{graph}{Black solid y=2sin(x),dashed y=2sin(x-2)+3,arrowed x,y axes,O; x ticks1 labels-4,-2,2,4; y ticks1 labels-4,-2,2,4,window-5..5,-5..5. Solid maxnear(1.57,2),min(-1.57,-2);dashed maxnear(-2.71,5),(3.57,5),min(.43,1);end arrows,equation labels.}
3. Complete the table to state the phase shift and vertical shift that will occur within the graph of the function.
\begin{tabular}{lll}
Function & Phase Shift & Vertical Shift \\
(y=3\cos x+2) & 0 & \answer{2} \\
(y=2\cos(x-1)+3) & \answer{1} & 3 \\
(y=\sin(x+\frac{3\pi}2)-1) & (3\pi/\answer{2}) & \answer{-1} \\
\end{tabular}
\te{Printed positive phase shift3pi/2 in last row; likely-3pi/2.}
enVision Algebra2. Teaching Resources. Copyright©Savvas Learning Company LLC.All Rights Reserved.
\end{te-addendum}
\begin{te-addendum}{Practice}{GeneralizeETP}
\textbf{7-4 Additional Practice. Translating Trigonometric Functions}
Name \underline{\hspace{2cm}}. enVision Algebra2. SavvasRealize.com.
Sketch the graph of the function.
1. (y=\cos(x+4))
\answer{Graph shown.}
% FIG 65 533 439 595 502
\placeholder{graph}{Magenta cos(x+4),arrowed x,y axes,O,unit xgrid labels-4,-2,2,4; ygrid1/2 labels-2,-1,1,2,window-5..5,-5/2..5/2. Max(-4,1),(2.28,1),min(-.86,-1),curvearrows.}
2. (y=\sin2(x-\frac\pi3))
\answer{Graph shown.}
% FIG 65 618 439 679 502
\placeholder{graph}{Magenta sin2(x-pi/3),arrowed x,y axes,O; x unitgrid labels-4,-2,2,4; yhalfunitgrid labels-2,-1,1,2,window-5..5,-5/2..5/2. Amplitude1 periodpi,max1.83+kpi,min.26+kpi,arrows.}
3. (y=3\sin(x-\frac\pi4)+2)
\answer{Graph shown.}
% FIG 65 704 439 766 506
\placeholder{graph}{Magenta3sin(x-pi/4)+2,arrowed x,y axes,O; unitgrid x,y; x labels-4,-2,2,4;ylabels2,4,6,window-5..5,-2..8. Max(2.36,5),min(-.79,-1),end arrows.}
Identify the amplitude,period,phase shift,vertical shift,and the maximum and minimum values of the function.
4. (y=\frac12\sin(x-\frac\pi2))
\answer{period:(2\pi);amplitude:(1/2);phase shift:(pi/2);vertical shift:0;max:0.5,min:-0.5}
5. (y=5\cos3(x+\pi)+6)
\answer{period:(2\pi/3);amplitude:5;phase shift:(-\pi);vertical shift:6;max:11,min:1}
6. The table shows the temperatures on different days of the year. Write a cosine model for the data. How does the midline value compare with the average of the12 temperatures?
\answer{(y\approx8.5\cos\frac{2\pi}{365}(x-196)+81.5);The average temp is close to the midline value.}
\begin{tabular}{lllllllllllll}
Day of the Year & 15 & 48 & 73 & 104 & 136 & 169 & 196 & 228 & 257 & 290 & 323 & 352 \\
Temperature(°F) & 76 & 73 & 75 & 79 & 82 & 87 & 90 & 89 & 88 & 87 & 83 & 79 \\
\end{tabular}
7. Use the data from Item6 to answer the following questions.
a. Write a sine function to model the weather station data.
\answer{Sample:(y=8.5\sin(\frac{2\pi}{365}(x-196)+\frac\pi2)+81.5)}
b. How do the cosine and sine models differ?
\answer{The sine model has a different phase shift compared to its parent function than the cosine model does.}
c. Use your sine model to estimate the temperature at the weather station on December31,day365.
\answer{about73°F}
enVision Algebra2.Teaching Resources.Copyright©Savvas Learning Company LLC.All Rights Reserved.
\end{te-addendum}
\begin{te-addendum}{Assess}{RtIExtend}
\textbf{7-4 Enrichment. Translating Trigonometric Functions}
Name \underline{\hspace{2cm}}. enVision Algebra2.SavvasRealize.com.
Nate and Phil model hiking trails through the woods as cosine or sine functions with the origin being a scenic outlook. Nate models the trails using sine functions and Phil models them using cosine functions. Match each of their equations to the corresponding trail to make sure both of their equations are accurate.
\begin{tabular}{lll}
Nate's Equations & Trails & Phil's Equations \\
(y=\pi\sin x+2) & Blue Trail & (y=4\cos(x-\frac\pi2)+\pi) \\
(y=4\sin x+2) & Red Trail & (y=\pi\cos(x-\frac\pi2)+2) \\
(y=4\sin x+\pi) & Green Trail & (y=4\cos(x-\frac\pi2)+2) \\
\end{tabular}
% FIG 65 860 454 1100 693
\placeholder{diagram}{Three-row matching chart Nate's Equations,Trails,Phil's Equations. Center three black sine graphs with arrowed x,y axes,O,x ticks1 labels-4,-2,2,4,window-5..5. Blue Trail4sin x+2,ywindow-3..7 with unitgrid labels-2,2,4,6,max6,min-2;Red Trail4sin x+pi,ywindow-2..8,unitgrid labels2,4,6,maxabout7.14,min-.86;Green Trailpi sin x+2,ywindow-4..6,labels-2,2,4,maxabout5.14,min-1.14. All curvearrows. Magenta matching lines connect lefttop toGreen,leftmiddletoBlue,leftbottomtoRed;righttoptoRed,rightmiddletoGreen,rightbottomtoBlue.}
\answer{Nate:(pi sin x+2)Green,(4sin x+2)Blue,(4sin x+pi)Red. Phil:(4cos(x-pi/2)+pi)Red,(pi cos(x-pi/2)+2)Green,(4cos(x-pi/2)+2)Blue.}
enVision Algebra2.Teaching Resources.Copyright©Savvas Learning Company LLC.All Rights Reserved.
\end{te-addendum}
\begin{te-addendum}{Assess}{ELLAddendum}
\textbf{7-4 Mathematical Literacy and Vocabulary. Translating Trigonometric Functions}
Name \underline{\hspace{2cm}}. enVision Algebra2.SavvasRealize.com.
Complete the vocabulary chart by matching the correct concept from the list.
a.the vertical translation
b.the distance for the graph to complete one full cycle
c.the average distance between the maximum and minimum of the graph
% FIG 65 167 1050 413 1118
\placeholder{graph}{Choices d,e,f:three black sine diagrams on unitgrids with arrowed x,y axes,O,xlabels-4,-2,2,4. d:solid4sin x,ywindow-5..5,labels-4,-2,2,4;dashed vertical doublearrow atxabout-1.6 spansy-4to4,labela. e:same solid4sin x plus dashed4cos x,with horizontal doublearrow near y2 fromx0 toaboutpi/2 labeledc;ywindow-5..5. f:solid4sin x,dashed4sin x+2,vertical doublearrownearxpi/2 fromy4to6,labeld;ywindow-4..6. All curvesend arrows.}
\begin{tabular}{lll}
Word or Word Phrase & Definition & Picture or Example \\
amplitude & The distance between the maximum and minimum of the graph. & 1.\answer{d.} \\
period & 2.\answer{b.} & Graph shown \\
vertical shift & 3.\answer{a.} & 4.\answer{f.} \\
midline & 5.\answer{c.} & Graph shown \\
phase shift & the horizontal translation & 6.\answer{e.} \\
\end{tabular}
\te{Printed amplitude definition is peak-to-trough distance; likely half this distance. Printed midline choice says average distance; likely average of maximum and minimum values.}
% FIG 65 342 1171 409 1239
\placeholder{graph}{Period example:black4sin x on unitgrid,arrowed x,y axes,O,xlabels-4,-2,2,4,ylabels-4,-2,2,4,window-5..5,-5..5. Dashed horizontal intervalarrow just belowxaxis spans-3.14to3.14,labeledb;end arrows on curve.}
% FIG 65 341 1258 410 1324
\placeholder{graph}{Midline example:black4sin x on unitgrid,arrowed x,y axes,O,xlabels-4,-2,2,4,ylabels-4,-2,2,4,window-5..5,-5..5. Dashed vertical doublearrow atxabout3.5 fromy-2to2,labelc near y1;end arrows.}
enVision Algebra2.Teaching Resources.Copyright©Savvas Learning Company LLC.All Rights Reserved.
\end{te-addendum}
\begin{te-addendum}{Assess}{GeneralizeETP}
\textbf{Digital Differentiated Resources}
The Reteach to Build Understanding,Additional Practice,and Enrichment activities are available as digital assignments. These activities are automatically assigned when students complete the lesson quiz online and are automatically scored.
\te{Tablet example screen: Exit. Solve the system of equations by graphing. (y=3x+4),(y=-2x+4). Use the graphing tool to graph the system. Click to enlarge graph. Question Help:Help Me Solve This;View an Example;Video;Textbook;Glossary;Math Tools;Print. Click the graph,choose a tool in the palette and follow the instructions to create your graph.1 part remaining.Clear All.Check Answer.Review progress.Question5 of14.Go.Back.Next.}
% FIG 65 583 985 1033 1303
\placeholder{illustration}{Black tablet with homebutton on left,white digital math page showing linear system3x+4,-2x+4 and blank coordinategrid at upperright. Grid x,y axes labeled-10,-8,...10,ticks1,origin0;QuestionHelpdropdown;blue footer with buttons and progressbar.}
\end{te-addendum}

\end{document}
